\gls{er} \cite{ll_exp_replay_2019} is a family of methods that mitigate \gls{cf} by maintaining a fixed-size replay buffer of past observations and interleaving them with current-task data during training. Two variants are used in this work, differing in their memory management strategies and training objectives: \gls{er-res}, which uses reservoir sampling, and \gls{er-ring}, which partitions the buffer evenly across tasks using a ring-buffer write policy.

\paragraph{Memory Management}
\textbf{\gls{er-res}} maintains an unstructured buffer $\mathcal{M}$ of capacity $K$ shared across all tasks. Samples are admitted via reservoir sampling \cite{ll_reservoir_sampling}: upon observing the $n$-th sample overall, it is written directly into the buffer if $|\mathcal{M}| < K$, otherwise it replaces a uniformly random existing entry with probability $K/n$:
%
\begin{equation}\label{eqn:reservoir}
    \mathcal{M} \leftarrow
    \begin{cases}
        \mathcal{M} \cup \{(x, y, t)\} & \text{if } |\mathcal{M}| < K \\
        \mathcal{M} \text{ with } \mathcal{M}[u] \leftarrow (x, y, t) & \text{if } |\mathcal{M}| = K
    \end{cases}
\end{equation}
where $u \sim \mathrm{Uniform}(0,\, n{-}1)$ and $n$ is the total number of samples seen so far.
%
This ensures that the buffer is an asymptotically uniform sample of all data seen so far, without requiring knowledge of the total number of future samples.

\textbf{\gls{er-ring}} partitions the buffer into $T$ fixed-size slots, one per task, with per-task capacity
%
\begin{equation}\label{eqn:er_ring_capacity}
    K_t = \left\lfloor \frac{K}{T} \right\rfloor + \mathbf{1}[t < r], \;\;\;\; r = K \bmod T
\end{equation}
%
where $T$ is the total number of tasks. Within each slot, samples are written sequentially via a circular pointer that wraps around when the slot is full.
This guarantees an equal per-task allocation throughout training, at the cost of requiring $T$ to be known in advance.

\paragraph{\gls{cl} Approach}

Both \gls{er-res} and \gls{er-ring} follow the same high-level replay loop. For each incoming batch from the current task, the algorithm samples a replay mini-batch from stored past examples, forms a joint training batch, and updates the shared model parameters by minimising cross-entropy on both current and replayed samples.

Let $B^\text{cur}$ denote the current-task batch and $B^\text{rep}$ the replay batch drawn from memory. The optimisation objective can be written as:
%
\begin{equation}\label{eqn:er_generic_loss}
    \mathcal{L}_\text{ER}
    = \frac{1}{|B^\text{cur}|}\sum_{(x,y)\in B^\text{cur}}
      \mathcal{L}_\text{CE}\!\left(f_\theta(x), y\right)
    + \frac{1}{|B^\text{rep}|}\sum_{(x,y,t)\in B^\text{rep}}
      \mathcal{L}_\text{CE}^{(t)}\!\left(f_\theta(x), y\right)
\end{equation}
%
where $f_\theta$ is the shared network, $\mathcal{L}_\text{CE}$ is standard cross-entropy, and $\mathcal{L}_\text{CE}^{(t)}$ denotes replay supervision with task-aware label/logit alignment.

In this implementation, the two variants differ in both memory structure and replay supervision:
\begin{itemize}
    \item \textbf{\gls{er-res}} uses a global reservoir buffer and replay classification over the full output head with global class labels.
    \item \textbf{\gls{er-ring}} uses per-task ring slots and task-conditioned replay alignment with task offsets and per-sample logit masking/gather.
\end{itemize}
Thus, \gls{er-res} directly trains cross-task class competition during replay, whereas \gls{er-ring} emphasises within-task consistency for replayed samples.

\paragraph{TIL and CIL Trade-offs}

The relative advantage of each memory structure depends on the evaluation protocol:
\begin{itemize}
    \item \textbf{\gls{til}:} \gls{er-ring} is often advantageous because fixed per-task allocation prevents old tasks from being underrepresented and task-conditioned replay matches the task-aware inference protocol. \gls{er-res} can still perform well, but its global reservoir may allocate memory unevenly when task sample counts differ.
    \item \textbf{\gls{cil}:} \gls{er-res} is often advantageous because global replay loss encourages direct discrimination among classes from different tasks. \gls{er-ring} benefits from balanced retention and explicit task metadata, but its task-conditioned replay objective can provide weaker pressure for cross-task class separation under strict single-head \gls{cil}.
\end{itemize}

Both approaches also involve practical capacity trade-offs. Reservoir replay fully utilises the shared buffer from early tasks onward, while fixed task-partitioning can leave future-task slots temporarily unused in early training. Conversely, fixed partitioning offers predictable per-task retention, which can reduce catastrophic forgetting when task data are imbalanced.

\paragraph{Inference}

Both variants use a standard linear classifier at inference time. In the \gls{til} setting, the task identity is used to select the appropriate logit slice or classifier head for the current task. In the \gls{cil} setting, the full logit vector is used and the class with the highest score across all seen classes is returned.